% ==== front matter: scope, conventions, the open problem (unnumbered) ====
\section*{Scope, conventions, and the open problem}
\addcontentsline{toc}{section}{Scope, conventions, and the open problem}

This is a self-contained mathematical reference for a sequential
\emph{active-learning} problem built on a Poisson--Gaussian-process model. A
latent function $\lat(x)$ over an input space is modeled by a Gaussian process
(GP) with an arc-cosine kernel over a structured prior covariance $\Cmat$. Each
input $x\in\Reals^{d}$ is a real-valued signal on a two-dimensional grid of $d$
grid points (e.g.\ a $108\times108$ grid, $d=11664$) confined to a bounded box.
Observations are nonnegative integer counts with a Poisson likelihood whose rate
is an exponential link of the latent, $\rate=\exp(\gain\lat+\latz)$. At each round
the active-learning rule chooses the next input so as to maximize the expected
information gain about $\lat$. Part~\ref{sec:diffusion} extends the loop from
\emph{selecting} the next input out of a fixed pool to \emph{synthesizing} it with
a guided diffusion model.

\paragraph{The open problem (the purpose of this document).} The arc-cosine kernel
is positively homogeneous in the input: its self-value $\Kf(x,x)=x\transpose\Cmat
x+\sigz^{2}$ grows with $\lVert x\rVert$. Under utility maximization the posterior
variance therefore grows like the \emph{square} of the input scale, the
acquisition utility diverges, and --- because inputs are confined to a bounded box
--- the optimizer drives them to the boundary of the box (saturation). Two remedies
are presented and analyzed: a normalized arc-cosine kernel and a saturating arc-sin
kernel (\S\ref{subsec:kernelfixes}). Neither is satisfactory in practice: the
normalized kernel is scale-invariant but fits poorly. It is an \textbf{open
problem} to find a kernel that simultaneously (i) keeps the acquisition utility
bounded on the bounded input domain and (ii) retains good predictive fit. The
divergence is proved in Part~\ref{sec:deeplayer} (\S\ref{subsec:divergence}); the
attempted fixes are \S\ref{subsec:kernelfixes}.

\paragraph{Conventions.}
\begin{itemize}[leftmargin=1.4em,itemsep=1pt,topsep=2pt]
  \item \textbf{One symbol, one concept} throughout; the full table is
        Appendix~\ref{app:notation}. A few symbols that collide across the
        literature are disambiguated there (the kernel smoothness $\smoo$ vs.\ the
        posterior correlation $\corr$; the prior width $\bwid$ vs.\ the diffusion
        noise schedule $\bsched$; the three entropies $\Hmarg,\Hnoise,\Hcond$).
  \item \textbf{Counts.} The count observation is written $\spk$; where the count
        index appears (as in the marginal entropy sum and its truncation
        $\rmax$) it is written $r$, denoting the same quantity, $\spk\equiv r$.
  \item \textbf{Scope of this rendering.} Content that is purely
        implementation-specific (automatic-vs-analytic differentiation machinery,
        the vector--Jacobian-product derivation, the low-level representation and
        online-update bookkeeping, and run/status details) has been omitted; the
        mathematics of the model, inference, utility, and the divergence problem is
        complete.
\end{itemize}

\paragraph{How to read this document.} The four parts are logically ordered.
Part~\ref{sec:genmodel}--\ref{sec:kernel} sets up the generative model and the
arc-cosine kernel with its structured prior. Part~II
(\ref{sec:varinf}--\ref{sec:training}) develops sparse variational inference, the
eigenspace representation of the posterior, and the coordinate-ascent training of
the model parameters. Part~III (\ref{sec:utility}--\ref{sec:deeplayer}) is the
active-learning core: the information utility and the proofs, subspace theory, and
numerical analysis beneath it --- including the divergence result and the attempted
kernel fixes that define the open problem. Part~IV (\ref{sec:diffusion}) is guided
diffusion for input synthesis. Cross-references connect the utility gradient in
Parts~III and IV back to the kernel input-gradient of Part~I, and the training
back to the ELBO of Part~II.
